\subsection{直和、多线性映射丛、对偶}

7.7.1. 假设 \(I_{-} = \emptyset\)。定义向量函子 \(\sigma\)，称为直和函子，规定 \(\sigma(\mathcal{V}) = \bigoplus_{i \in I} V_i\) 且 \(\sigma(f) = \bigoplus_{i \in I} f_i\)。若 \(\mathcal{M} = (M^i)_{i \in I}\) 是以 B 为基底的向量丛族，则向量丛 \(\sigma(\mathcal{M})\) 称为 \(M^i\) 的直和，记为 \(\bigoplus_{i \in I} M^i\)。对每个 \(b \in B\)，在 \(b\) 处的 \(\bigoplus_{i \in I} M^i\) 的纤维是各个 \(M^i\) 在 \(b\) 处的纤维的直和。设 U 是 B 的一个开子集，且 \(s_i \in \mathscr{S}_{M^i}^r(U)\)（\(i \in I\)）。则映射 \(b \mapsto \sum_{i \in I} s_i(b)\) 是一个截面，记为 \(\sum_{i} s_i\)，其类为 \(C^r\)，属于 \(M = \bigoplus_{i \in I} M^i\)；并且映射 \((s_i)_{i \in I} \mapsto \sum_{i} s_i\) 是一个 \(\mathscr{C}^r(U)\)-模同构，从 \(\bigoplus_{i} \mathscr{S}_{M^i}^r(U)\) 到 \(\mathscr{S}_M^r(U)\)。

\(\bigoplus_{i\in I}M^{i}\) 的底层流形与纤维积 \(\prod_{B}M^{i}\) 相同。

将 \(pr_{i}\) 记为从 \(\bigoplus_{i\in I} M^{i}\) 到 \(M^{i}\) 的向量丛态射；它在每个纤维 \(\bigoplus_{i\in I} M_{b}^{i}\) 上是第 i 个投影。类似地，定义典范嵌入 \(j_{i}\)，它将 \(M^{i}\) 嵌入 \(\bigoplus M^{i}\)。

设 \(f\) 是从 B 到流形 \(\mathbf{B}'\) 的态射；设 H 是另一个有限集，\(\mathcal{N} = (\mathrm{N}^{h})_{h\in \mathbf{H}}\) 是以 \(\mathbf{B}'\) 为基底的向量丛族。映射 \(u\mapsto \bigl(\mathrm{pr}_{h}\circ u\circ j_{i}\bigr)_{(h,i)\in \mathbf{H}\times \mathbf{I}}\) 是 \(f\)-态射集，从 \(\bigoplus_{i\in \mathbf{I}}\mathbf{M}^i\) 到 \(\bigoplus_{h\in \mathbf{H}}\mathbf{N}^h\)，到矩阵集 \((u_{h,i})_{(h,i)\in \mathbf{H}\times \mathbf{I}}\) 的双射，其中 \(u_{h,i}\) 是一个 \(f\)-态射，从 \(\mathbf{M}^i\) 到 \(\mathbf{N}^h\)。

若 I = \{1, 2\}，则序列

\[
0 \to \mathbf {M}^{1} \xrightarrow{j_{1}} \mathbf {M}^{1} \oplus \mathbf {M}^{2} \stackrel{\mathrm{pr}_{2}}{\to} \mathbf {M}^{2} \to 0
\]

是直正合的。

反之，设 M 是以 B 为基底的向量丛，M' 是 M 的子向量丛。假设流形 B 是仿紧的，且满足以下两个条件之一：

\begin{enumerate}
\def\labelenumi{(\roman{enumi})}
\item
  K 不同于 R 或 C；
\item
  \(\mathbf{K} = \mathbf{R}, r \neq \omega\)，且流形 \(\mathbf{B}\) 具有类为 \(C^r\) 的单位分解（5.3.6）。
\end{enumerate}

则存在 M 的一个子向量丛 M''，使得 M 可与直和 M' ⊕ M'' 认同。

7.7.2. 假设 \(I_{+}=\{0\}\) 且 \(I_{-}=\{1,2,\ldots,d\}\)。定义类型为 I 的向量函子 \(\eta_{d}\)，其类为 \(C^{r}\)，规定 \(\eta_{d}(\mathcal{V})=\mathcal{L}(V_{1},\ldots,V_{d};V_{0})\)，并规定 \(\eta_{d}(\mathbf{f})(u)=f_{0}\circ u\circ(f_{1}\times\cdots\times f_{d})\)，其中 \(u\in\eta_{d}(\mathcal{V})\)。若 \(\mathcal{M}=(M_{i})_{i\in I}\) 是以 B 为基底的向量丛族，则将向量丛 \(\eta_{d}(\mathcal{M})\) 记为 \(\mathcal{L}(M_{1},\ldots,M_{d};M_{0})\)。

设 \(u\) 是从 \(\mathbf{M}_{1} \times_{\mathbf{B}} \cdots \times_{\mathbf{B}} \mathbf{M}_{d}\) 到 \(\mathbf{M}_{0}\) 的多线性态射。映射 \(\hat{u}: b \mapsto u_{b}\) 于是是 \(\mathcal{L}(\mathbf{M}_{1}, \ldots, \mathbf{M}_{d}; \mathbf{M}_{0})\) 的一个截面，而映射 \(u \mapsto \hat{u}\) 是双射。

7.7.3. 保留 7.7.2 的记号，并进一步假设 \(d=1\)。向量丛 \(\mathcal{L}(\mathbf{M}_{1};\mathbf{M}_{0})\) 称为从 \(\mathbf{M}_{1}\) 到 \(\mathbf{M}_{0}\) 的同态丛。它的截面与从 \(\mathbf{M}_{1}\) 到 \(\mathbf{M}_{0}\) 的 B-态射相对应。

此外，若 \(M_{0}\) 是平凡丛 \(K_{B}\)，则向量丛 \(\mathcal{L}(\mathbf{M}_{1};\mathbf{K}_{\mathbf{B}})\) 称为 \(M = M_{1}\) 的对偶，记为 \(M'\)：纤维 \((\mathbf{M}')_{b}\) 是 M 的纤维 \(M_{b}\) 上的连续线性形式空间，其中点 \(b \in B\)。

若 \(s\)（分别 \(t\)）是 M（分别 M'）在 B 的一个开集 U 上的截面，则映射 \(b \mapsto (b, \langle s(b), t(b) \rangle)\) 是一个截面，记为 \(\langle s, t \rangle\)；这个截面属于平凡丛 \(\mathbf{K}_{\mathbf{B}}\)。¹

